\begin{theorembox}{Theorem (Heine--Borel Theorem for $\mathbb{R}$)}
\begin{theorem}[Heine--Borel Theorem for $\mathbb{R}$]
\label{thm:heine-borel}
Let $K\subseteq\mathbb{R}$.  Then $K$ is compact if and only if $K$ is closed
and bounded.  Explicitly,
\[
  K\text{ is compact}
  \Longleftrightarrow
  \leqft[
  \mathbb{R}\setminus K\text{ is open}
  \right]
  \land
  \leqft[
  (\exists M>0)(\forall x\in K)(|x|\leq M)
  \right].
\]
\hyperref[prf:heine-borel]{Go to proof.}
\end{theorem}
\end{theorembox}
\begin{remark*}[Standard quantified statement]
\[
K \text{ compact} \Longleftrightarrow (K \text{ closed} \land K \text{ bounded}).
\]
\end{remark*}
\begin{remark*}[Interpretation]
Heine--Borel is the bridge into Analysis~I: it certifies that the closed bounded sets --- the ones on which the Extreme Value and uniform-continuity theorems live --- are exactly the compact ones. It is stated here as the equivalence \emph{on $\mathbb{R}$}; the open-cover side is the definition that later carries to the torus, while ``closed and bounded'' does not.
\end{remark*}
\begin{remark*}[Exposition]
With Heine--Borel in hand the chain is complete: number construction $\to$ embedded number systems $\to$ the real line as an ordered metric space $\to$ Analysis~I. Completeness, introduced as the last rung of the embedding ladder, is exactly the property that makes the line compact in the large, and compactness is the hypothesis Analysis~I leans on.
\end{remark*}
\begin{dependencies}
\begin{itemize}
  \item \hyperref[thm:compact-implies-closed-bounded]{Compact Subsets of $\mathbb{R}$ Are Closed and Bounded}
  \item \hyperref[thm:closed-bounded-interval-compact]{Closed Bounded Intervals Are Compact}
\end{itemize}
\end{dependencies}
